% ============================================================================
% Fichier  : partie4/P04_C16_incidents.tex
% Titre    : Gérer les Incidents de Sécurité
% Auteur   : MINKA MI NGUIDJOÏ Thierry Emmanuel
% Version  : 1.0 -- novembre 2026
% ============================================================================
% Sources : ISO/IEC 27035-1:2023 ; NIST SP 800-61 Rev. 3 (2025). Registre
% des incidents de LiZbiZ : cas fictif.
% ============================================================================

\chapter{Gérer les Incidents de Sécurité}
\label{chap:incidents}
\chaptermark{Gérer les incidents}

\epigraph{%
    « Un incident mal classé est un incident mal traité. »%
}{Principe du cours}

\niveauChapitre{Intermédiaire}{Chapitres~\ref{chap:traiter}
et~\ref{chap:rapport}~; cours SEC3266, chapitre sur les opérations.}

\begin{boxobjectifs}
\begin{itemize}
    \item Distinguer événement, incident et crise.
    \item Dérouler les phases de la gestion des incidents selon
          l'ISO/IEC~27035.
    \item Classer un incident et l'escalader selon une grille.
    \item Organiser l'équipe de réponse et ses liens avec la direction.
    \item Mesurer la gestion des incidents et en tirer des enseignements
          pour le registre.
\end{itemize}
\end{boxobjectifs}

\bigskip

% ============================================================================
\section*{Question directrice : à partir de quand un incident devient-il l'affaire du directeur général~?}
% ============================================================================

Le 5~mars à 22~h~05, le technicien d'astreinte voit que l'application ne
répond plus. Il redémarre les serveurs et rentre chez lui~: pour lui,
c'est une panne. Le lendemain à 7~h~30, c'est une crise. Entre les deux,
aucune grille ne disait qui prévenir, ni à partir de quand. Ce chapitre
traite de la gestion des incidents du point de vue de l'organisation.

\medskip
\textbf{Plan du chapitre.}
\begin{enumerate}
    \item Événement, incident, crise (section~\ref{sec:c16-definitions}).
    \item Les phases de l'ISO/IEC 27035 (section~\ref{sec:c16-phases}).
    \item Classer et escalader (section~\ref{sec:c16-classer}).
    \item L'équipe de réponse (section~\ref{sec:c16-equipe}).
    \item Mesurer et apprendre (section~\ref{sec:c16-mesurer}).
\end{enumerate}

\separateur

% ============================================================================
\section{Événement, Incident, Crise}
\label{sec:c16-definitions}
% ============================================================================

\begin{boxdef}
\termecle{Événement de sécurité}~: occurrence qui indique une possible
violation de la sécurité ou une défaillance des mesures.
\termecle{Incident de sécurité}~: un ou plusieurs événements indésirables
qui ont une probabilité significative de compromettre l'activité ou de
menacer la sécurité de l'information~\cite{iso27035_2023}. Une
\termecle{crise} est une situation qui dépasse les capacités de la gestion
courante et exige des décisions de la direction (Chapitre~\ref{chap:crise}).
\end{boxdef}

\begin{tableaucours}{L{3.0cm}L{4.6cm}L{4.6cm}}{%
    \tableauHeader{Niveau} &
    \tableauHeader{Exemple} &
    \tableauHeader{Qui gère}
}
Événement & Une alerte d'antivirus bloquée & Exploitation, sans escalade \\
Incident & Un poste infecté, isolé & Équipe de réponse \\
Incident majeur & Un compte à privilèges compromis & Équipe de réponse,
RSSI, direction informée \\
Crise & Serveurs chiffrés sur deux sites & Cellule de crise présidée par le
DG \\
\end{tableaucours}
\vspace{-0.8em}
\begin{center}{\small\itshape Tableau~16.1. De l'événement à la crise}\end{center}

% ============================================================================
\section{Les Phases de l'ISO/IEC 27035}
\label{sec:c16-phases}
% ============================================================================

L'\norme{ISO/IEC~27035-1:2023} décrit cinq phases de gestion des
incidents~\cite{iso27035_2023}. Le NIST SP~800-61 Rev.~3 (2025) les relie
aux fonctions du NIST CSF~2.0~\cite{nist80061r3}.

\begin{tableaucours}{L{3.4cm}L{8.8cm}}{%
    \tableauHeader{Phase} &
    \tableauHeader{Contenu}
}
Planifier et préparer & Politique, équipe, procédures, outils, grille de
gravité, contacts, exercices \\
Détecter et signaler & Recevoir les alertes et les signalements, les
enregistrer \\
Évaluer et décider & Qualifier l'événement, décider s'il s'agit d'un
incident, le classer, l'escalader \\
Répondre & Confiner, éradiquer, rétablir, préserver les preuves,
communiquer \\
Tirer les enseignements & Analyser les causes, améliorer les mesures, le
registre et les procédures \\
\end{tableaucours}
\vspace{-0.8em}
\begin{center}{\small\itshape Tableau~16.2. Les cinq phases de l'ISO/IEC 27035-1:2023}\end{center}

\begin{boxalert}
\textbf{Violation typique.} Sauter l'évaluation (« on redémarre et on
verra »).

\smallskip
\textbf{Règle~:} tout événement inhabituel sur un système critique est
évalué selon la grille avant d'agir~; en cas de doute, on escalade. Une
escalade inutile coûte un appel~; une escalade manquée a coûté, en mars,
une nuit et des preuves.
\end{boxalert}

% ============================================================================
\section{Classer et Escalader}
\label{sec:c16-classer}
% ============================================================================

\begin{tableaucours}{L{2.0cm}L{5.4cm}L{4.8cm}}{%
    \tableauHeader{Gravité} &
    \tableauHeader{Critère} &
    \tableauHeader{Escalade}
}
4 · Critique & Activité arrêtée sur plusieurs sites~; données sensibles
exfiltrées~; atteinte publique & DG et cellule de crise, sous 1~heure \\
3 · Haute & Compte à privilèges compromis~; un site arrêté & RSSI et
direction, sous 4~heures \\
2 · Moyenne & Poste ou serveur isolé touché, sans propagation & Équipe de
réponse, dans la journée \\
1 · Basse & Tentative bloquée, signalement & Traitement courant \\
\end{tableaucours}
\vspace{-0.8em}
\begin{center}{\small\itshape Tableau~16.3. Grille de gravité et d'escalade de LiZbiZ (convention du cours)}\end{center}

\begin{boxscenario}{Opération GAROUA. Ce que la grille aurait changé}
\textbf{Constatation.} À 22~h~05 le 5~mars, l'application de facturation
ne répond plus à Douala et à Garoua, et un message de rançon s'affiche sur
un serveur.

\medskip
\textbf{Analyse} (confiance élevée). Selon la grille, deux sites arrêtés
et un message de rançon font un incident de gravité~4~: escalade au DG
sous une heure, soit avant 23~h~05. La cellule de crise aurait été réunie
dans la nuit, au lieu de 10~h~30 le lendemain.
\end{boxscenario}

% ============================================================================
\section{L'Équipe de Réponse}
\label{sec:c16-equipe}
% ============================================================================

\begin{tableaucours}{L{3.8cm}L{8.4cm}}{%
    \tableauHeader{Rôle} &
    \tableauHeader{Responsabilité}
}
Responsable de l'incident & Coordonne, décide dans son périmètre, escalade \\
Analystes & Qualifient, confinent, éradiquent, collectent les preuves \\
Prestataire & Agit sur les systèmes qu'il administre, selon le contrat \\
Juridique & Preuves, plainte, obligations envers les tiers \\
Communication & Messages internes et externes, validés \\
Direction & Arbitre au-delà du périmètre, décide en cas de crise \\
\end{tableaucours}
\vspace{-0.8em}
\begin{center}{\small\itshape Tableau~16.4. Les rôles de la réponse à incident}\end{center}

\begin{boxprocedure}
\textbf{Ouvrir un incident chez LiZbiZ}
\begin{enumerate}
    \item Enregistrer l'incident (numéro, heure, source, description).
    \item Classer selon la grille~; en cas de doute, prendre le niveau
          supérieur.
    \item Escalader selon le délai de la grille, par téléphone, puis par
          écrit.
    \item Désigner le responsable de l'incident et ouvrir une main
          courante.
    \item Préserver avant d'agir~: ne pas éteindre, isoler, collecter.
    \item Reclasser à chaque point de situation.
\end{enumerate}
\end{boxprocedure}

% ============================================================================
\section{Mesurer et Apprendre}
\label{sec:c16-mesurer}
% ============================================================================

\begin{tableaucours}{L{4.8cm}L{7.4cm}}{%
    \tableauHeader{Indicateur} &
    \tableauHeader{Définition}
}
Délai de détection & Du début de l'incident à sa détection \\
Délai d'escalade & De la détection à l'escalade au niveau de la grille \\
Délai de confinement & De la détection au confinement \\
Délai de rétablissement & De la détection au retour à la normale \\
Part des incidents au registre & Incidents dont le scénario figurait au
registre des risques (Chapitre~\ref{chap:auditer-risques}) \\
\end{tableaucours}
\vspace{-0.8em}
\begin{center}{\small\itshape Tableau~16.5. Indicateurs de la gestion des incidents}\end{center}

\begin{boxobjet}{R-02 · Leçons de l'incident de mars}
\textbf{Scénario~:} usurpation de l'accès distant du prestataire.
\textbf{Délais constatés~:} détection, 80~h~18~min~; escalade au DG,
15~minutes après la détection~; confinement (coupure des liaisons),
3~h~40~; rétablissement de la facturation, 11~jours. \textbf{Leçons~:}
grille d'escalade, astreinte formée, guide de réponse au rançongiciel.
\textbf{Propriétaire~:} M.~KOUNOU. \textbf{Indicateur~:} délai d'escalade
des incidents de gravité~3 et~4.
\end{boxobjet}

\begin{boxmetier}{Responsable de la continuité}
\textbf{Relier incident et continuité.} Le responsable de la continuité
reçoit toute escalade de gravité~3 ou~4~: c'est lui qui juge s'il faut
basculer en mode dégradé ou activer le plan de continuité.
\end{boxmetier}

\separateur

\begin{boxresume}
\textbf{Ce qu'il faut retenir de ce chapitre~:}
\begin{itemize}
    \item \textbf{Événement, incident, crise}~: trois niveaux, trois
          gestionnaires.
    \item \textbf{L'ISO/IEC~27035}~: préparer, détecter, évaluer, répondre,
          apprendre.
    \item \textbf{Une grille} classe et fixe le délai d'escalade~; dans le
          doute, on escalade.
    \item \textbf{L'équipe de réponse} associe technique, prestataire,
          juridique, communication et direction.
    \item \textbf{On mesure} détection, escalade, confinement,
          rétablissement, et l'on nourrit le registre.
\end{itemize}
\end{boxresume}

\begin{boxinfo}
\textbf{Points de contrôle de l'élève.} Je sais~: distinguer événement,
incident et crise~; nommer les cinq phases~; classer un incident et
l'escalader~; décrire les rôles de la réponse~; définir les délais qui la
mesurent.
\end{boxinfo}

\bigskip

\begin{boxexo}[title={\ding{45}\quad Exercice~16.1. Quel niveau~? \niveaubase}]
Classez en événement, incident, incident majeur ou crise~: une alerte
d'hameçonnage bloquée~; l'agence de Garoua sans réseau pendant six
heures~; la publication de données clients.
\end{boxexo}

\begin{boxexo}[title={\ding{45}\quad Exercice~16.2. Les phases \niveaubase}]
Dans quelle phase de l'ISO/IEC~27035 se rangent~: l'exercice annuel~; la
plainte~; la mise à jour du registre après l'incident~?
\end{boxexo}

\begin{boxexo}[title={\ding{45}\quad Exercice~16.3. Le délai \niveaubase}]
D'après la grille, avant quelle heure le DG devait-il être prévenu le
5~mars~?
\end{boxexo}

\begin{boxexo}[title={\ding{45}\quad Exercice~16.4. L'astreinte \niveaumoyen}]
Rédigez la fiche réflexe de l'astreinte de LiZbiZ pour un serveur qui ne
répond plus la nuit.
\end{boxexo}

\begin{boxexo}[title={\ding{45}\quad Exercice~16.5. Reclasser \niveaumoyen}]
Un incident ouvert en gravité~2 (un serveur touché) révèle, deux heures
plus tard, l'usage d'un compte d'administrateur du domaine. Que faites-vous~?
\end{boxexo}

\begin{boxexo}[title={\ding{45}\quad Exercice~16.6. Le prestataire \niveaumoyen}]
Quelles obligations le contrat de SAHEL INFOGÉRANCE doit-il prévoir pour la
gestion des incidents~?
\end{boxexo}

\begin{boxexo}[title={\ding{45}\quad Exercice~16.7. Les délais \niveaumoyen}]
Calculez, pour l'incident de mars, le délai entre le chiffrement
(5~mars, 21~h~40) et la détection (6~mars, 7~h~30). Pourquoi est-il plus
parlant que les 80~heures pour la direction~?
\end{boxexo}

\begin{boxexo}[title={\ding{45}\quad Exercice~16.8. Le registre \niveaumoyen}]
Pourquoi mesurer la part des incidents dont le scénario figurait au
registre~? Que signifie un taux faible~?
\end{boxexo}

\begin{boxexo}[title={\ding{45}\quad Exercice~16.9. Une procédure \niveauexpert}]
Rédigez la procédure de gestion des incidents de LiZbiZ selon les cinq
phases, en deux pages.
\end{boxexo}

\begin{boxexo}[title={\ding{45}\quad Exercice~16.10. Une grille bancaire \niveauexpert}]
Adaptez la grille du Tableau~16.3 à la BANQUE SANAGA, en tenant compte des
clients, des autorités de supervision et des systèmes de paiement.
\end{boxexo}

\bigskip

\begin{boxinfo}
\textbf{Pour aller plus loin.}
\begin{itemize}
    \item ISO/IEC~27035-1~\cite{iso27035_2023}.
    \item NIST SP~800-61 Rev.~3~\cite{nist80061r3}.
\end{itemize}
\end{boxinfo}

\begin{boxinfo}
\textbf{Transition.} Pour décider d'un mode dégradé, il faut savoir ce que
chaque activité peut supporter. Le Chapitre~\ref{chap:bia} présente
l'analyse d'impact sur l'activité.
\end{boxinfo}
